% Final-presentation work-in-progress deck.
% No animations and no paper screenshots; all equations are typeset in LaTeX.
\documentclass[aspectratio=169,11pt]{beamer}
\usepackage[utf8]{inputenc}
\usepackage[T1]{fontenc}
\usepackage[english]{babel}
\usepackage{amsmath,amssymb}
\usepackage{booktabs}
\usepackage{xcolor}
\usetheme{default}
\usecolortheme{seahorse}
\setbeamertemplate{navigation symbols}{}
\setbeamertemplate{footline}[frame number]
\setbeamerfont{frametitle}{size=\large}
\definecolor{ink}{RGB}{32,48,71}
\definecolor{accent}{RGB}{30,113,120}
\setbeamercolor{frametitle}{fg=ink}
\setbeamercolor{structure}{fg=accent}

\newcommand{\kb}{k_B}
\newcommand{\km}{k_M}
\newcommand{\Rb}{R_B}
\newcommand{\Rm}{R_M}

\title{Credit Quantity, Credit Prices, and Future Output}
\subtitle{A model of heterogeneous financing access\\Final presentation --- work in progress}
\author{Alvaro Marcelo Chavez Unyen}
\institute{Universidad del Pacifico}
\date{October 2026}

\begin{document}

\begin{frame}[plain]
  \titlepage
  \begin{center}\small
    \href{https://github.com/amchavezu/ai-project}{\texttt{github.com/amchavezu/ai-project}}
  \end{center}
\end{frame}

\begin{frame}{1 \textperiodcentered{} Question and conditional answer}
\textbf{Question.} Why can credit quantity be more informative for one sector,
while the price of credit is more informative for another?

\medskip
\textbf{Answer.} The relevant indicator depends on which financing margin
determines the marginal unit of capital:
\begin{itemize}
  \item at a binding cheap-bank cap, quantity determines feasible investment;
  \item with outside finance at the margin, its common price determines desired investment.
\end{itemize}

\medskip
\alert{The thesis evidence is preliminary: the model explains a qualitative
short-versus-longer-horizon pattern, not exact forecast peaks.}
\end{frame}

\begin{frame}{2 \textperiodcentered{} Agent, timing, and financing order}
At date $t$, a firm with internal funds $n$ chooses capital $k$, bank credit
$b$, and outside finance $m$. Capital produces at $t+1$.
\[
\max_{k,b,m}\ \beta F(k)-\Rb b-\Rm m
\quad\text{s.t.}\quad
k=n+b+m,\quad 0\le b\le\bar b,\quad m\ge0.
\]
\[
F'>0,\quad F''<0,\quad 0<\Rb\le\Rm<\beta F'(0).
\]
Bank finance is cheaper but capped; the outside-finance wedge captures access,
issuance, and information costs.
\end{frame}

\begin{frame}{3 \textperiodcentered{} Piecewise optimum}
Let $\bar k=n+\bar b$, and define the two unconstrained targets by
\[
\beta F'(\kb)=\Rb,\qquad \beta F'(\km)=\Rm,
\qquad \km\le\kb.
\]
Then
\[
k^*=\begin{cases}
n, & \kb\le n,\\
\kb, & n<\kb<\bar k,\\
\bar k, & \km\le\bar k\le\kb,\\
\km, & \bar k<\km.
\end{cases}
\]
The empirical mechanism compares the third line (quantity regime) with the
fourth line (price regime).
\end{frame}

\begin{frame}{4 \textperiodcentered{} Main result: two transmission margins}
\small
\begin{tabular}{p{0.27\linewidth}p{0.31\linewidth}p{0.31\linewidth}}
\toprule
 & Binding bank cap & Outside finance at margin \\
\midrule
Strict regime & $\km<\bar k<\kb$ & $\bar k<\km$ \\
Capital & $k^*=n+\bar b$ & $k^*=\km(z,\tau)$ \\
Bank-line effect & $\partial k^*/\partial\bar b=1$ & $\partial k^*/\partial\bar b=0$ \\
Spread effect & $\partial k^*/\partial z=0$ locally &
$\partial k^*/\partial z=1/[\beta F''(\km)]<0$ \\
Economic signal & feasible investment pipeline & marginal cost of total finance \\
\bottomrule
\end{tabular}

\medskip
Zero effects are local: sufficiently large shocks can cross regime thresholds.
\end{frame}

\begin{frame}{5 \textperiodcentered{} From firm heterogeneity to sectors}
Let $\lambda_s$ be the share of sector $s$ in the quantity regime:
\[
Y_{s,t+1}=\lambda_s F(n+\bar b_s)
+(1-\lambda_s)F\!\left(\km(z_s,\tau_s)\right).
\]
If $0\le\lambda_T<\lambda_N\le1$ and conditional primitives are common,
\[
\frac{\partial Y_N}{\partial\bar b}>
\frac{\partial Y_T}{\partial\bar b}\ge0,
\qquad
\left|\frac{\partial Y_N}{\partial z}\right|<
\left|\frac{\partial Y_T}{\partial z}\right|.
\]
This is a conditional composition result. Whether $\lambda_N>\lambda_T$ in
Peru is an empirical assumption to test, not a theorem.
\end{frame}

\begin{frame}{6 \textperiodcentered{} Mechanism and empirical checks}
\begin{enumerate}
  \item \textbf{Bank dependence:} quantity predictive content should rise with
  exhausted credit lines or dependence on relationship banks.
  \item \textbf{Alternative funding:} spread predictive content should rise
  with access to bonds, foreign credit, or other external finance.
  \item \textbf{Substitution:} a larger bank line should raise total investment
  at the cap, but mainly replace outside finance for market-access firms.
\end{enumerate}

\medskip
Project gestation and adjustment costs can make the price channel appear over a
longer horizon. The present static proposition does not identify those dynamic
parameters.
\end{frame}

\begin{frame}{7 \textperiodcentered{} Checks and numerical illustration}
For $\beta=.95$, $a=2$, $c=.5$, $n=.4$, $\bar b=.5$, $\Rb=.5$, $\Rm=1.4$:
\[
\bar k=.900,\qquad \kb=2.947,\qquad \km=1.053.
\]
\begin{center}
\small
\begin{tabular}{lcc}
\toprule
Experiment & Quantity regime & Price regime \\
\midrule
$\bar b$ rises by $.10$ & capital $\uparrow$ to $1.000$ & capital unchanged \\
$z$ rises by $.02$ & capital unchanged & capital $\downarrow$ to $1.011$ \\
\bottomrule
\end{tabular}
\end{center}
\texttt{code/verify.py} independently checks the FOCs, regime inequalities,
comparative statics, and all reported values; it fails if a check is false.
\end{frame}

\begin{frame}{8 \textperiodcentered{} Formalization and audit trail}
\begin{itemize}
  \item The hand appendix derives the policy and both propositions step by step.
  \item The current Lean 4 prototype proves the quadratic core and sectoral
  inequalities with no \texttt{sorry}.
  \item The required AppliedModelingLib-generated folder, pinned paper commit,
  and \texttt{check --fast} output remain pending for the final deadline.
  \item The paper appendix maps every numbered result to its Lean declaration.
  \item \texttt{prompts.md} preserves the raw AI exchanges and the paper records
  accepted, corrected, and rejected claims.
\end{itemize}
\end{frame}

\begin{frame}{9 \textperiodcentered{} Limitations and next decisions}
\textbf{Known limits}
\begin{itemize}
  \item credit ceilings and sector composition are reduced-form primitives;
  \item the measured loan spread may not equal the marginal outside-finance cost;
  \item the forecasting relation is not a causal identification result.
\end{itemize}

\medskip
\textbf{Feedback sought}
\begin{enumerate}
  \item Is the two-tranche model the right minimum mechanism for the thesis?
  \item Should the next extension endogenize the ceiling or model dynamics first?
  \item Which Peruvian proxy best measures alternative-finance access by sector?
\end{enumerate}
\end{frame}

\end{document}
